\documentclass{book}
\usepackage{ctex}
\usepackage[cmex10]{amsmath}
\usepackage{xcolor}
\usepackage{mathtools}
\usepackage{booktabs}
\DeclareRobustCommand{\bluebold}[1]{\textbf{\textcolor[rgb]{0.31,0.63,0.86}{#1}}}
\usepackage[
    pdfpagemode=UseOutlines,
    bookmarks=true,
    bookmarksopen=true,
    colorlinks=true,
    linkcolor=black,
    citecolor=black,
    urlcolor=black
]{hyperref}
\usepackage{bookmark}

\raggedbottom

\begin{document}
\tableofcontents
\clearpage


流体流动基本方程知识点

\begin{table}[htbp]
\centering
\begin{tabular}{c c} 
\toprule
知识点 & point\\ [0.5ex] 
\midrule
 1 & 连续性方程 \\
 2 & 伯努利方程 \\
\bottomrule
\end{tabular}
\caption{Knowledge Points Table of Basic Equations of Fluid Flow.}
\label{table:data}
\end{table}

\chapter{连续性方程}

稳态流动，管路无分支和汇合，质量守恒

\begin{equation}
q_\text{m1} = q_\text{m2}
\end{equation}

\begin{equation}
\rho_1 q_\text{v1} = \rho_2 q_\text{v2}
\end{equation}

不可压缩流体：

\begin{equation}
\rho_1 = \rho_2 = \rho
\end{equation}

\begin{equation}
q_\text{v1} = q_\text{v2}
\end{equation}

不可压缩流体在\bluebold{圆形管路}中流动

\begin{equation}
\begin{rcases}
\displaystyle A = \frac{\pi}{4} d^2 \\[4pt]
\displaystyle q_\text{v} = uA \\[4pt]
\displaystyle q_\text{v1} = q_\text{v2}
\end{rcases}
\quad
\Rightarrow
\quad
\displaystyle u_1 d_1^2 = u_2 d_2^2
\quad \Rightarrow \quad
\frac{u_1}{u_2} = \left(\frac{d_2}{d_1}\right)^2
\label{eq:continuity-set}
\end{equation}

稳态流动系统，流量一定，圆形管道中各截面流速仅与直径有关

\section*{tips}

\bluebold{计算相关会与机械能衡算式、流动阻力计算等融合在一起作为一道大题考}

\chapter{伯努利方程}

\section{伯努利方程}

衡算范围：管内壁面、1-1'与 2-2' 间截面

衡算基准：1kg 流体主流区

基准水平面：0-0' 平面

对于做稳态流动的不可压缩流体，机械能衡算

\begin{equation}
gz_1 + \frac{u_1^2}{2} + \frac{p_1}{\rho} + W_\mathrm{e} = gz_2 + \frac{u_2^2}{2} + \frac{p_2}{\rho} + h_\mathrm{f}
\qquad \mathrm{J/kg}
\label{eq:bernoulli}
\end{equation}

符号说明（单位质量，J/kg）

\vspace{6pt}
$gz_1$、$gz_2$：位能，相对值，基于同一基准水平面

\vspace{6pt}
\(\dfrac{p_1}{\rho}\)、\(\dfrac{p_2}{\rho}\)：静压能，基于同一压强基准，绝对压强或表压

\vspace{6pt}
\(\dfrac{u_1^2}{2}\)、\(\dfrac{u_2^2}{2}\)：动能，和速度有关

\vspace{6pt}
$h_\mathrm{f}$：阻力损失，流体具有黏性导致的

\vspace{6pt}
有效功 $W_\mathrm{e}$：输送设备对 1kg 流体做的有效功

\vspace{6pt}
有效功率 $N_\mathrm{e}$：单位时间输送设备做的有效功率，J/s 或 W

\begin{equation}
N_\mathrm{e} = q_\text{m} W_\mathrm{e} = \rho q_\text{v} W_\mathrm{e}
\qquad \mathrm{J/s} \quad \text{or} \quad \mathrm{W}
\end{equation}

\section*{另一种表示方法}

化工原理常用 $h_\mathrm{f}$（单位质量能量损失）

\begin{equation}
gz_1 + \frac{u_1^2}{2} + \frac{p_1}{\rho} + W_\mathrm{e} = gz_2 + \frac{u_2^2}{2} + \frac{p_2}{\rho} + h_\mathrm{f}
\end{equation}

符号说明（单位质量，J/kg）

\vspace{6pt}
gz：位能

\vspace{6pt}
\(\dfrac{u^2}{2}\)：动能

\vspace{6pt}
\(\dfrac{p}{\rho}\)：静压能

\vspace{6pt}
$W_\mathrm{e}$：外加功（泵对流体做功）

\vspace{6pt}
$h_\mathrm{f}$：流动阻力损失

\section{应用伯努利方程解题要点}

\bluebold{确定衡算范围}：做出流动示意图，标明流体流动方向，选定上、下游截面，明确衡算范围

\bluebold{衡算截面}：截面与流动方向垂直，截面间流体连续，待求的未知量应在截面间或截面上。

\bluebold{基准面}：任意水平面，一般选位置较低的截面中心线所在的水平面


\bluebold{压强}：表示方法一致，绝对压强或表压强，表压强更方便

\bluebold{单位一致}：采用国际单位

\end{document}
